\documentclass[11pt]{article}
\usepackage[T1]{fontenc}
\usepackage{textcomp}
\usepackage[margin=0.75in]{geometry}
\usepackage{amsmath,amssymb,amsthm}
\usepackage{array,booktabs}
\usepackage{hyperref}
\setlength{\parindent}{0pt}
\newtheorem{theorem}{Theorem}
\theoremstyle{definition}
\newtheorem{definition}{Definition}
\title{Flow Proof Artifact: prop-14-equal-chords-equidistant}
\author{Generated by \texttt{flow doc proof}}
\date{Flow Proof Book}
\begin{document}
\maketitle
In equal circles equal straight lines are equally distant from the centers, and conversely.\\[0.5em]
\textbf{Source.} euclid — Elements, Book III, Proposition 14\\[0.5em]
\bigskip
\noindent{\large \textbf{Derived fact 1.} \textit{In equal circles equal straight lines are equally distant from the centers, and conversely} \hfill the Euclidean plane \cdot Euclid Book III \cdot Proposition 14: equal chords are equally distant from the center}\par
\smallskip\noindent\textit{Coordinate: the Euclidean plane \cdot Euclid Book III \cdot Proposition 14: equal chords are equally distant from the center}\par
\smallskip\noindent\textit{Source: euclid — Elements, Book III, Proposition 14}\par
\smallskip\noindent\textit{Built on: proposition 2: a perpendicular from the center bisects the chord, for Euclid Book III on the Euclidean plane}\par
\medskip
\noindent\textbf{Goal.} In equal circles equal straight lines are equally distant from the centers, and conversely\par
\noindent$\displaystyle \text{equal chords equally distant}$\par
\medskip
\noindent
\begin{tabular}{@{} >{\raggedright\arraybackslash}p{0.06\textwidth} >{\raggedright\arraybackslash}p{0.47\textwidth} >{\raggedleft\arraybackslash}p{0.41\textwidth} @{}}
\textbf{\#} & \textbf{Proof} & \textbf{Mathematics} \\
\midrule
\textbf{1.} & We prove that proposition 14: equal chords are equally distant from the center for Euclid Book III the Euclidean plane. & \\[0.45em]
\textbf{2.} & Let AB = chord\_first. & \\[0.45em]
\textbf{3.} & Let DE = chord\_second. & \\[0.45em]
\textbf{4.} & Let d1 = distance\_to\_center. & \\[0.45em]
\textbf{5.} & From step 2, step 3, and step 4, we can deduce that d1 equals d2. & $\displaystyle d1 = d2$ \\[0.45em]
\textbf{6.} & We invoke the derived fact governing Euclid Book III the Euclidean plane: proposition 2: a perpendicular from the center bisects the chord, for Euclid Book III on the Euclidean plane. & \\[0.45em]
\textbf{7.} & From step 6, this implies equal chords equally distant. Hence proven. & $\displaystyle \text{equal chords equally distant}$ \\[0.45em]
\bottomrule
\end{tabular}
\label{thm:Geometry-Euclid Book III-proposition_14_equal_chords_are_equally_distant_from_the_center}
\par\medskip\hrule
\end{document}
